\vfill\pagebreak
\section{Blocks}
\label{sec:control:blocks}

Every construction of this chapter works on \textit{blocks}, so they come first.
A block is a sequence of instructions between braces, \token{\{} and \token{\}}.
The body of \token{main} is one. A block can also stand anywhere a value is
expected, and then it does two things.

\begin{itemize}
\item It \textbf{groups} instructions. They run in order, from the opening
  brace to the closing one, and the variables declared inside exist only until
  the closing brace. Past that point their names mean nothing.
\item It \textbf{has a value}: the value of its last expression, provided that
  expression is not followed by a semicolon.
\end{itemize}

In the next listing, the block on lines~5 to~8 computes a price with a tax of
one fifth. The variable \token{tax} is a step of the calculation and nothing
more, so it lives inside the block. The block's value is \token{price + tax},
written last and without a semicolon, and that is what \token{total} receives.

\begin{lstlisting}[style=coloredverbatim, caption=A block used as a value, label=lst:control:block_value]
use std::io;

fn main() {
  let price = read!i32(ask-> "Price: ");
  let total = {
    let tax = price / 5;
    price + tax
  };
  println("Total with tax: ", total);
}
\end{lstlisting}
\vspace{-10pt}%
\begin{lstlisting}[style=bashVerb, escapechar=@+]
@+\textcolor{teal!80}{alice@dev:\mtilde}@+$ ./main
Price: 100
Total with tax: 120
\end{lstlisting}

The two properties each have a way to go wrong. Once the block is closed,
\token{tax} no longer exists, and a line~9 that tried to print it would stop the
compilation with ``undefined symbol tax''.

The other mistake is to end the last expression with a semicolon. The semicolon
turns an expression into an instruction whose value is thrown away, so the block
no longer has a value to give. Its type is then \token{void}, the type of
``nothing'', and no variable can hold nothing.

\begin{lstlisting}[style=coloredverbatimError, escapechar=@, caption=A block whose last expression ends with a semicolon]
use std::io;

fn main() {
  let price = read!i32(ask-> "Price: ");
  let @\hb{total}@ = {
    let tax = price / 5;
    price + tax;
  };
  println("Total with tax: ", total);
}
\end{lstlisting}

The compiler reports that it ``cannot declare a variable of type void''. The
line it points at is the declaration of \token{total}, but the fault is the
semicolon on line~7.

A block does not have to give a value. It can also stand where an instruction is
expected, between two other instructions, and then it only groups. That is
useful on its own: a variable needed for a few lines, declared in such a block,
disappears at the closing brace, and the rest of the function cannot use it by
mistake. Blocks are used most, though, in the constructions that follow, where
each branch of a decision and each repetition of a loop is a block.
